\chapter{GIỚI THIỆU ĐỀ TÀI}

\section{Đặt vấn đề}

Trong bối cảnh toàn cầu hóa, ngày càng nhiều tổ chức doanh nghiệp vận hành đồng thời tại
nhiều quốc gia với ngôn ngữ khác nhau. Các doanh nghiệp hoạt động tại khu vực châu Á --- nơi
hội tụ nhiều ngôn ngữ như Tiếng Nhật, Tiếng Trung, Tiếng Việt, Tiếng Thái, Tiếng Bengali,
Tiếng Hindi và Tiếng Indonesia --- phải xử lý khối lượng lớn tài liệu kỹ thuật cần được dịch
nhất quán và chính xác: báo cáo kiểm tra chất lượng, hướng dẫn vận hành, thông số kỹ thuật,
biên bản xác nhận.

Thực trạng cho thấy bốn thách thức cốt lõi:

\begin{enumerate}
  \item \textbf{Thiếu nhất quán thuật ngữ chuyên ngành.} Các công cụ dịch phổ thông như
    Google~Translate hay DeepL không hỗ trợ bộ từ điển thuật ngữ riêng của tổ chức. Cùng
    một thuật ngữ kỹ thuật (tên sản phẩm, tên quy trình, ký hiệu tiêu chuẩn) có thể được
    dịch khác nhau giữa các tài liệu, gây nhầm lẫn trong vận hành và giao tiếp nội bộ.

  \item \textbf{Hỗ trợ định dạng tài liệu hạn chế.} Tài liệu doanh nghiệp tồn tại dưới nhiều
    định dạng: văn bản DOCX, bảng tính XLSX, thuyết trình PPTX, và tệp PDF --- bao gồm cả
    PDF có thể chỉnh sửa (native text) và PDF dạng scan (ảnh). Phần lớn các công cụ hiện có
    không giữ nguyên định dạng, bố cục và kiểu chữ sau khi dịch.

  \item \textbf{Thiếu quy trình quản lý thuật ngữ tập trung.} Không có hệ thống lưu trữ và
    quản lý từ điển thuật ngữ dùng chung, dẫn đến mỗi bộ phận hoặc cá nhân tự duy trì
    danh sách thuật ngữ riêng, không đồng bộ và khó kiểm soát chất lượng.

  \item \textbf{Không xử lý được PDF scan.} Nhiều tài liệu kỹ thuật lịch sử tồn tại dưới
    dạng PDF scan (ảnh chụp hoặc scan từ bản in), không thể trích xuất nội dung văn bản để
    dịch bằng các công cụ thông thường.
\end{enumerate}

Những thách thức trên đặt ra nhu cầu về một hệ thống dịch tài liệu được thiết kế cho môi
trường doanh nghiệp đa ngôn ngữ, tích hợp quản lý thuật ngữ chuyên ngành và hỗ trợ nhiều
định dạng tài liệu.

\section{Mục tiêu và phạm vi đề tài}

\subsection*{Mục tiêu}

Đồ án đặt mục tiêu xây dựng hệ thống web dịch tài liệu đa ngôn ngữ dành cho doanh nghiệp với
các tính năng cốt lõi sau:

\begin{itemize}
  \item Hỗ trợ dịch sang \textbf{10 ngôn ngữ}: Tiếng Việt (vi), Tiếng Nhật (ja), Tiếng Trung
    giản thể (zh-CN), Tiếng Trung phồn thể (zh-TW), Tiếng Anh (en), Tiếng Thái (th), Tiếng
    Bengali (bn), Tiếng Hindi (hi), Tiếng Indonesia (id) và Tiếng Oriya (or).

  \item Hỗ trợ \textbf{6 định dạng tài liệu}: DOCX, XLSX, PPTX, PDF native-text, PDF scan,
    CSV/XLS --- với bảo toàn định dạng sau khi dịch.

  \item Xây dựng \textbf{hệ thống quản lý thuật ngữ ba cấp}: Thư viện chung (Common Library)
    quản lý bởi Admin/Library Keeper, Thư viện cá nhân (Private Library) cho từng người dùng,
    và hàng chờ đề xuất (Suggestion Queue) với cơ chế duyệt cộng đồng.

  \item Giao diện web đơn trang (SPA) trực quan, hỗ trợ dịch nhiều ngôn ngữ đích đồng thời
    trong một yêu cầu.

  \item Hệ thống quản trị toàn diện: quản lý tài khoản người dùng theo vai trò, thống kê sử
    dụng thuật ngữ, theo dõi lịch sử dịch.

  \item Xác thực người dùng qua AWS Cognito SSO, triển khai trên hạ tầng đám mây AWS.
\end{itemize}

\subsection*{Phạm vi đề tài}

\begin{itemize}
  \item \textbf{Trong phạm vi:} Người dùng trong tổ chức doanh nghiệp đa ngôn ngữ; dịch tài
    liệu từ các định dạng trên sang bất kỳ ngôn ngữ nào trong 10 ngôn ngữ hỗ trợ; quản lý
    từ điển thuật ngữ chuyên ngành theo cơ chế đóng góp cộng đồng.

  \item \textbf{Ngoài phạm vi:} Dịch thời gian thực (real-time streaming), ứng dụng di động
    (mobile app), tổng hợp giọng nói (text-to-speech), dịch nội dung website hay hệ thống
    quản trị nội dung (CMS).
\end{itemize}

\section{Định hướng giải pháp}

Hệ thống được thiết kế theo kiến trúc ba tầng với các thành phần sau:

\begin{itemize}
  \item \textbf{Tầng giao diện (Frontend):} Ứng dụng web đơn trang React~18 + Vite~6 phục vụ
    qua Nginx, sử dụng Tailwind CSS v4 và Ant Design v5 cho giao diện người dùng.

  \item \textbf{Tầng logic nghiệp vụ (Backend):} Django~6.0 với Django REST Framework~3.16
    cung cấp REST~API; PostgreSQL lưu trữ dữ liệu người dùng, lịch sử dịch và từ điển thuật
    ngữ; AWS~S3 lưu trữ tệp gốc và tệp đã dịch.

  \item \textbf{Dịch vụ dịch thuật:} Google Cloud Translation~API~v3 với 44~cặp glossary tùy
    chỉnh lưu trên Google Cloud Storage; Azure Document Intelligence cho OCR tài liệu scan;
    ConvertAPI cho chuyển đổi định dạng (PDF→DOCX, XLS→XLSX).

  \item \textbf{Xác thực:} Amazon Cognito OIDC (SSO môi trường production) + JWT
    (djangorestframework-simplejwt) cho môi trường phát triển.
\end{itemize}

\section{Bố cục đồ án}

Đồ án được tổ chức thành sáu chương:

\begin{itemize}
  \item \textbf{Chương 1:} Trình bày bối cảnh, đặt vấn đề, mục tiêu và định hướng giải pháp của
    đồ án.

  \item \textbf{Chương 2:} Khảo sát hiện trạng các công cụ dịch thuật hiện có; phân tích và
    đặc tả yêu cầu chức năng và phi chức năng thông qua biểu đồ use case và bảng đặc tả.

  \item \textbf{Chương 3:} Trình bày nền tảng lý thuyết về dịch máy thần kinh, glossary, OCR,
    và các công nghệ, thư viện được lựa chọn cùng lý do lựa chọn.

  \item \textbf{Chương 4:} Trình bày thiết kế kiến trúc, thiết kế chi tiết (giao diện, lớp,
    cơ sở dữ liệu), kết quả xây dựng, kiểm thử và triển khai hệ thống.

  \item \textbf{Chương 5:} Trình bày ba giải pháp và đóng góp kỹ thuật nổi bật: pipeline phân
    loại PDF, hệ thống thuật ngữ ba cấp và cơ chế inject context qua thread-local.

  \item \textbf{Chương 6:} Kết luận về kết quả đạt được, hạn chế còn tồn tại và định hướng
    phát triển trong tương lai.
\end{itemize}
